\chapter{Tlamantli Tlamaniliztli}
\begin{abstract}

Inin amatl tlatlaçotla in tlamaniliztli para in tlapetlanilotl
tēpoztli ātlacualoyan, tlen mochihua ica in tlatlaminaliztli
ihuan in tlatlapaliztli tlen āyōtl.
Inin tlamaniliztli OpenLaval quipalehuia in tlamaniliztli
ica ome īxiptla: tlen se tlamantli (symmetric) ihuan tlen
ome‑ixiptla (asymmetric), in tlen quipiya in Prandtl–Meyer
tlapalehuiliztli para īpan ihuan īcxitlan in ātlacualoyan.

Ipan in ome‑ixiptla tlamantli, in īpan ihuan īcxitlan
mopalehuiya ica ahmo se tlamantli tlatlapaliztli,
huan yeh inin quichihua in R$^\ast$ ahmo se,
in tlatlapaliztli ahmo se, ihuan in tlatzintla‑tlacuilolli
ahmo se. Inin amatl quipohua in tlamantli tlen
OpenLaval quichihua axcan.

\end{abstract}

\section{Tlayecantli}

In tlamaniliztli tlen tlapetlanilotl tēpoztli
mopehua ica in tlacuilolli tlen quixpolōa in tlatlaminaliztli
tlen mochihua īpan in tlayacantli.
Mochi tlen monequi:

\begin{itemize}
  \item inlet Mach number $M_{\mathrm{in}}$
  \item outlet Mach number $M_{\mathrm{out}}$
  \item specific heat ratio $\gamma$
  \item inlet flow angle $\beta_{\mathrm{in}}$
  \item Prandtl–Meyer angles para īpan ihuan īcxitlan
\end{itemize}

OpenLaval quipalehuia se‑ixiptla tlamantli
(ica se par $(\nu_l, \nu_u)$)
ihuan ome‑ixiptla tlamantli
(ica nāhui tlapalehuiliztli):



\[
\nu_{l,\mathrm{lower}},\quad
\nu_{u,\mathrm{lower}},\quad
\nu_{l,\mathrm{upper}},\quad
\nu_{u,\mathrm{upper}}.
\]



Inin tlapalehuiliztli monequi quipiya in tlanahuatilli:



\[
\nu_{\min} \le \nu_l < \nu_u \le \nu_{\max},
\]



ken quipiya $\nu_{\min} = \min(\nu(M_{\mathrm{in}}), \nu(M_{\mathrm{out}}))$ ihuan 
$\nu_{\max} = \max(\nu(M_{\mathrm{in}}), \nu(M_{\mathrm{out}}))$.

\section{Tlamaniliztli Tlatlamantli}

In tlamaniliztli mochihua inin:

\begin{enumerate}
  \item Tlatlapaliztli tlen quixpolōa in tlatlaminaliztli
        para īpan ihuan īcxitlan.
  \item Tlachihualiztli in tlatzintla‑tlacuilolli (arc segments)
        para cada īxiptla (R$^\ast$ ahmo se).
  \item Tlachihualiztli in īpan tlacuilolli ica in
        $\beta_{\mathrm{in}}$ ihuan quitzontequi para īcxitlan.
  \item Tlatzontequiliztli in huehcatl (blade height)
        ihuan tlatzontequiliztli in mochi coordenadas.
  \item Tlayacantli‑tlatequi (leading‑edge trimming)
        ica \texttt{edge\_delta} ihuan \texttt{edge\_offset}.
\end{enumerate}

Ipan in ome‑ixiptla tlamantli, in 1–3 mochihua
para cada īxiptla.

\section{Ome‑Ixiptla Prandtl–Meyer Tlapalehuiliztli}

Ipan in ome‑ixiptla tlamantli, in tlacuilo quipiya
nāhui tlapalehuiliztli:



\[
(\nu_{l,\mathrm{lower}}, \nu_{u,\mathrm{lower}}), \qquad
(\nu_{l,\mathrm{upper}}, \nu_{u,\mathrm{upper}}).
\]



Inin quipiya in tlatlapaliztli tlen mochihua īpan
ihuan īcxitlan, tlen monequi quipiya $\nu_{\min} \le \nu_l < \nu_u \le \nu_{\max}$ īpan cada īxiptla.
Yeh inin quichihua in tlatlapaliztli ahmo se.

OpenLaval quichihua tlatlanahuatilli para mochi
tlapalehuiliztli.

\section{Tlatlaminaliztli‑Quixpolōlli (Shock‑Cancelling Curves)}

In tlatlapaliztli mochihua ica in tlanahuatilli:



\[
\phi(R^\ast, \theta) = \mathrm{const.}
\]



Ipan in ome‑ixiptla tlamantli, in tlayecantli
ahmo se, ihuan in R$^\ast$ ahmo se.

\section{Tlatzintla‑Tlacuilolli (Arc Segments)}

Para quichihua se tlatzontequiliztli,
cada īxiptla quipiya se arc segment.
Ipan in ome‑ixiptla tlamantli,


\[
R^\ast_{\mathrm{upper}} \ne R^\ast_{\mathrm{lower}},
\]


huan in radius ahmo se.

\section{Tlayecantli‑Tlacuilolli (Straight‑Line Region)}

In īpan tlacuilolli mopalehuiya ica in
$\beta_{\mathrm{in}}$ ihuan quitzontequi para īcxitlan.
Ipan in ome‑ixiptla tlamantli, in tlatzontequiliztli
ahmo se.

\section{Tlayacantli‑Tlatequi (Leading‑Edge Trimming)}

OpenLaval quichihua in tlayacantli‑tlatequi
ica \texttt{edge\_delta} ihuan \texttt{edge\_offset}.
Ipan in ome‑ixiptla tlamantli, in tlatequi‑tlacuilolli
ahmo se.

\section{Program Usage}

In config amatl quipiya:

\begin{itemize}
  \item symmetric PM: \texttt{vl}, \texttt{vu}
  \item asymmetric PM:
        \texttt{vl\_lower}, \texttt{vl\_upper},
        \texttt{vu\_lower}, \texttt{vu\_upper}
  \item inlet angle: \texttt{beta\_in}
  \item leading‑edge: \texttt{edge\_delta}, \texttt{edge\_offset}
  \item interpolation: \texttt{num\_points}
\end{itemize}

In CLI quipiya \texttt{--asymmetric} para mochi
tlapalehuiliztli.
